\documentclass[11pt]{article}
\usepackage[margin=1in]{geometry}
\usepackage{amsmath,amssymb,amsthm}
\usepackage{graphicx}
\usepackage[hidelinks]{hyperref}

\newtheorem{theorem}{Theorem}
\newtheorem{lemma}{Lemma}
\newtheorem{definition}{Definition}
\newtheorem{remark}{Remark}

\newcommand{\dens}{\operatorname{dens}}
\newcommand{\cf}{\operatorname{cf}}
\newcommand{\lin}{\operatorname{lin}}
\newcommand{\omegaone}{\omega_1}

\title{A partial positive answer for overcomplete sets in ordinal interval spaces}
\author{Agent lane 16}
\date{2026-06-23}

\begin{document}
\maketitle

\noindent\textbf{Status.}
Candidate partial result, likely valid. This packet proves a ZFC positive
subcase of Question 48 from Koszmider's arXiv:2006.00806, not the full
\(\xi<\omega_2\) problem.

\section{Source question}

The source paper is:

\begin{quote}
P. Koszmider, \emph{On the existence of overcomplete sets in some classical
nonseparable Banach spaces}, arXiv:2006.00806.
\end{quote}

On page 25, Question 48 asks:

\begin{quote}
Can one prove in ZFC that the Banach spaces \(C([0,\xi])\) for all ordinals
\(\xi<\omega_2\) admit overcomplete sets?
\end{quote}

\begin{center}
\includegraphics[width=0.92\linewidth]{figures/open_problem_crop.png}
\end{center}

\section{Definitions and source tool}

\begin{definition}
Let \(X\) be a Banach space. A set \(D\subseteq X\) is \emph{overcomplete}
if \(|D|=\dens(X)\) and every subset \(A\subseteq D\) with
\(|A|=|D|\) has dense linear span in \(X\).
\end{definition}

We use the following criterion from the source paper, Theorem 10 there. For
reference we state only the form needed here.

\begin{lemma}[Koszmider's criterion]\label{lem:criterion}
Let \(X\) be a Banach space admitting a linearly dense family
\(\{x_\alpha:\alpha<\omega_1\}\). Suppose that there is a norm closed subspace
\(Y\subseteq X^*\) of finite codimension such that, for every \(y^*\in Y\),
\[
   \{\alpha<\omega_1: y^*(x_\alpha)\neq 0\}
\]
is countable. Then some finite-codimensional subspace of \(X\) admits an
overcomplete set. More precisely, there is \(k<\omega\) such that every
codimension-\(k\) subspace of \(X\) admits an overcomplete set.
\end{lemma}

The source paper applies this with \(k\leq 1\) to \(C([0,\omega_1])\), and then
uses the standard fact that, for the scattered compacta under consideration,
finite-codimensional subspaces of \(C(K)\) are isomorphic to \(C(K)\). We use
the same last step.

\section{Idea of the proof}

Koszmider's proof for \(C([0,\omega_1])\) produces a dense family whose
interaction with measures is countable after one endpoint mass is removed.
For ordinals below \(\omega_1\cdot\omega\), the compact interval \([0,\xi]\)
breaks into finitely many clopen ordinal blocks, each of type at most
\(\omega_1\). On each block the same countable-support criterion is available:
countable blocks are separable, and \(\omega_1\)-blocks are exactly the source
case. A finite \(\ell_\infty\)-sum preserves the criterion because a finite
union of countable exceptional sets is countable and the codimensions add only
finitely.

\section{The partial theorem}

\begin{theorem}\label{thm:main}
In ZFC, for every ordinal \(\xi<\omega_1\cdot\omega\), the Banach space
\(C([0,\xi])\) admits an overcomplete set.
\end{theorem}

\begin{proof}
If \(\xi<\omega_1\), then \(C([0,\xi])\) is separable, and the existence of an
overcomplete sequence is the classical theorem of Klee, quoted as Theorem 4 in
the source paper. Thus assume \(\omega_1\leq \xi<\omega_1\cdot\omega\).

Choose \(m<\omega\) such that \(\xi<\omega_1\cdot m\). The ordinal interval
\([0,\xi]\) is the finite union of clopen intervals
\[
   K_0,\ldots,K_{m-1},
\]
where empty final pieces are omitted, and each \(K_j\) is order-homeomorphic
to an ordinal interval \([0,\eta_j]\) with \(\eta_j\leq\omega_1\). Concretely,
one cuts at the finitely many points \(\omega_1,\omega_1\cdot2,\ldots\) below
\(\xi\); intervals of the form \((\gamma,\delta]\) in an ordinal compactum are
clopen. Therefore
\[
   C([0,\xi]) \cong \bigoplus_{j<r}^{\ell_\infty} C(K_j)
\]
for some finite \(r\).

We first check that each summand \(C(K_j)\) satisfies the hypothesis of
Lemma~\ref{lem:criterion}, allowing a finite-codimensional subspace of its
dual. If \(\eta_j<\omega_1\), then \(C(K_j)\) is separable; choose a countable
linearly dense set \(D_j\), and take \(Y_j=C(K_j)^*\). For each \(y^*\in Y_j\),
the set of \(d\in D_j\) with \(y^*(d)\neq0\) is countable simply because
\(D_j\) is countable.

If \(\eta_j=\omega_1\), then \(K_j\) is homeomorphic to \([0,\omega_1]\), and
the source proof of Theorem 11(2) gives exactly the needed data: a linearly
dense set \(D_j\subseteq C(K_j)\) of cardinal \(\omega_1\) and a codimension-one
subspace \(Y_j\subseteq C(K_j)^*\) such that every \(y^*\in Y_j\) is nonzero on
at most countably many members of \(D_j\).

Now pass to the finite \(\ell_\infty\)-sum
\[
   X=\bigoplus_{j<r}^{\ell_\infty} C(K_j).
\]
Embed each \(D_j\) into the \(j\)-th coordinate of \(X\), and let
\[
   D=\bigcup_{j<r} \iota_j(D_j).
\]
The set \(D\) is linearly dense in \(X\). Since at least one block has density
\(\omega_1\) in the present nonseparable case, \(|D|=\omega_1\).

The dual is the finite \(\ell_1\)-sum
\[
   X^*=\bigoplus_{j<r}^{\ell_1} C(K_j)^*.
\]
Set
\[
   Y=\bigoplus_{j<r}^{\ell_1} Y_j.
\]
Then \(Y\) has finite codimension in \(X^*\), since each \(Y_j\) has finite
codimension and there are only finitely many summands. If
\(y^*=(y_j^*)_{j<r}\in Y\), then \(y^*\) is nonzero on the coordinate copy
\(\iota_j(D_j)\) only at the countably many points where \(y_j^*\) is nonzero
on \(D_j\). Hence \(y^*\) is nonzero on only a finite union of countable sets
inside \(D\), therefore on only countably many members of \(D\).

Thus \(X\cong C([0,\xi])\) satisfies Lemma~\ref{lem:criterion}. Consequently a
finite-codimensional subspace of \(C([0,\xi])\) admits an overcomplete set.
Finally, \([0,\xi]\) is an infinite compact ordinal interval with a nontrivial
convergent sequence. As used in the source proof of Theorem 11(2), this implies
that finite-codimensional subspaces of \(C([0,\xi])\) are isomorphic to
\(C([0,\xi])\). Since admitting an overcomplete set is preserved by Banach-space
isomorphism, \(C([0,\xi])\) itself admits an overcomplete set.
\end{proof}

\section{Limitations and novelty check}

This does not solve Question 48 in full. The argument depends on having only
finitely many \(\omega_1\)-sized ordinal blocks. It therefore reaches
\(\xi<\omega_1\cdot\omega\), but it does not address ordinals such as
\(\omega_1\cdot\omega\), \(\omega_1^2\), or arbitrary \(\xi<\omega_2\), where
one would need either a countable-block version of Koszmider's criterion or a
new selection argument.

Local run indexes were searched for arXiv:2006.00806, overcomplete-set
keywords, \(C([0,\xi])\), Grothendieck/scattered variants, and related ordinal
phrases. The only close hit was a separate candidate full packet for
arXiv:1912.08690 concerning weakly discrete overcomplete sets in
\(C([0,\omega_1])\); it does not cover the present ordinal-interval subcase.
Web searches for the exact Question 48 wording and close variants did not find
a later paper resolving this subcase or the full \(\xi<\omega_2\) problem.

\section{Human review recommendation}

Review as a likely valid partial result. The main checks are:
\begin{enumerate}
\item the finite clopen decomposition of \([0,\xi]\) for
      \(\xi<\omega_1\cdot\omega\);
\item the finite-sum preservation of Lemma~\ref{lem:criterion};
\item the final use of the source paper's finite-codimensional
      \(C(K)\cong\) hyperplane fact for compact ordinal intervals with
      nontrivial convergent sequences.
\end{enumerate}

\begin{thebibliography}{9}

\bibitem{Koszmider}
P. Koszmider,
\emph{On the existence of overcomplete sets in some classical nonseparable
Banach spaces},
arXiv:2006.00806.

\bibitem{RussoSomaglia}
T. Russo and J. Somaglia,
\emph{Overcomplete sets in non-separable Banach spaces},
arXiv:1912.08690.

\end{thebibliography}

\end{document}
